% Shared body of the anonymised (main_anon.tex) and named (main_biorxiv.tex) versions.
% Every number is a macro from numbers.tex (scripts/91_paper_numbers.py).

\begin{abstract}
World models that predict the consequences of actions in a learned latent space have been
built for robots and games; whether the same recipe can model how a population of human cells
evolves under a genetic perturbation is open. We train an action-conditioned latent world model
of genome-scale live-cell RNAi time-lapse microscopy: a frozen self-supervised video encoder
(V-JEPA~2) turns two-day movies into latent trajectories, and a small transformer predicts the
next latent state given the past and the knocked-down gene. We ask three pre-registered
questions, each with a decision gate fixed before the test data were read: whether conditioning
on the gene improves forecasts for held-out plates beyond an action-agnostic model and a
replicate-memorisation baseline; whether gene representations from independent modalities
(Perturb-seq, Cell Painting, protein language models, interaction networks) transfer the model
to genes never seen in training; and whether searching over gene actions recovers the genes
that produce a target phenotype or trajectory. On test, the gene-conditioned model beat both
comparators ($p=\SignalTestAgP$), but a model trained with randomly permuted reagent labels did
equally well, the effect did not replicate on an independent reagent set, zero-shot transfer
passed for one of five gene representations only marginally, and planning was not
distinguishable from nearest-gene retrieval. We report this as a negative result and describe
the controls that exposed it.
\end{abstract}

\section{Introduction}

A world model predicts how the state of an environment changes under an action and can then be
used to plan, by searching for actions whose predicted consequences reach a goal
\citep{ha2018worldmodels,hafner2025dreamer}. Recent self-supervised video encoders make this
possible without pixels: V-JEPA~2 learns latent video representations from internet-scale video
and, with a small action-conditioned predictor trained on its frozen latents, plans robot
manipulation in latent space \citep{assran2025vjepa2}. Cell biology has its own version of the
problem. A genetic perturbation is an action, a time-lapse movie of a cell population is a
trajectory, and the questions a world model answers (what will happen under this knockdown, and
which knockdown produces this phenotype) are central to functional genomics and to the
``virtual cell'' programme \citep{bunne2024virtualcell,xing2026worldmodel}.

Existing virtual-cell world models predict transcriptomic responses
\citep{qiu2026chreode,miao2026vclmu,chuai2026alphacell}, and image-based perturbation models
generate static images of perturbed cells \citep{palma2025impa,wang2025morphdiff,
zhang2025cellflux,zhang2026cellfluxv2} or 3D morphologies \citep{naidoo2025form}. Live-cell
video, where a perturbation's phenotype unfolds over hours, has been modelled generatively
without perturbation conditioning \citep{holmberg2025videodiffusion}. We found no prior model
that conditions latent dynamics of live-cell video on the perturbed gene and evaluates transfer
to genes held out entirely (Section~\ref{sec:related}).

We build such a model on the MitoCheck screen, a genome-wide RNAi screen in which every
knockdown was imaged for two days by time-lapse microscopy \citep{neumann2010mitocheck}, and we
design its evaluation around three lessons from perturbation modelling. First, simple baselines
are hard to beat \citep{ahlmanneltze2025baselines}, so every model must beat a
memorisation baseline that replays the training replicates of the same reagent. Second, plate
and batch effects can dominate learned image representations, so the metric is computed after a
per-plate normalisation to negative controls and the plate identity carried by the latent is
measured. Third, a test set evaluated more than once stops being a test set, so the design,
gates, primary metrics and evaluation lists were frozen before training, and the test splits are
evaluated once.

Our contributions are: (i) an action-conditioned latent world model of genome-scale live-cell
RNAi video, trained on \NWellsStageTwoEmbedded{} wells beyond a first stage of
\NMoviesStageOne{} movies; (ii) a pre-registered evaluation with three decision gates
(prediction, zero-shot transfer and planning), memorisation and kNN baselines, trained
shuffled-encoding controls and per-plate normalisation; (iii) a quantification of the batch
structure carried by frozen video-encoder latents of microscopy; and (iv)
a negative result: post hoc controls show that a model trained
with permuted reagent labels matches the gene-conditioned model, so the pre-registered
prediction gain does not establish gene-conditioned dynamics.

\section{Related work}\label{sec:related}

\paragraph{World models.} Latent world models learn a transition model in a learned
representation and plan by searching over actions \citep{ha2018worldmodels,hafner2025dreamer}.
V-JEPA~2 \citep{assran2025vjepa2} separates the two: a self-supervised video encoder is
pre-trained without actions, then frozen, and an action-conditioned predictor (V-JEPA~2-AC) is
trained on its latents with teacher forcing and a rollout loss. We follow this recipe with gene
knockdowns as actions.

\paragraph{Virtual cells and perturbation prediction.} Models of perturbation responses have
mostly predicted transcriptomes \citep{roohani2024gears}, and simple linear baselines remain
competitive with deep models for unseen perturbations \citep{ahlmanneltze2025baselines}.
Recent ``cell world models'' treat perturbations as actions on a latent cell state, again for
expression data \citep{qiu2026chreode,miao2026vclmu,chuai2026alphacell}. Perspectives argue for
models that couple modalities, including imaging \citep{bunne2024virtualcell,xing2026worldmodel}.

\paragraph{Image-based perturbation models.} Generative models predict how cells look after a
genetic or chemical perturbation from static images: IMPA by style transfer
\citep{palma2025impa}, MorphDiff by diffusion conditioned on L1000 expression
\citep{wang2025morphdiff}, CellFlux and CellFluxV2 by flow matching between control and perturbed
populations \citep{zhang2025cellflux,zhang2026cellfluxv2}, and FORM for 3D morphology
\citep{naidoo2025form}. Video diffusion has been adapted to time-lapse microscopy, conditioned on
phenotype descriptions or a first frame rather than a perturbation
\citep{holmberg2025videodiffusion}.

\paragraph{Live-cell RNAi screens.} MitoCheck imaged a genome-wide siRNA library in HeLa cells
expressing histone~2B--GFP and scored mitotic phenotypes per cell with supervised classifiers
\citep{neumann2010mitocheck,held2010cellcognition}; later work found phenotypes without labels
using hidden Markov models on the class time courses \citep{failmezger2013unsupervised}. We use
the movies through the Image Data Resource and the BioImage Archive
\citep{williams2017idr,hartley2022bia}.

\paragraph{Novelty search.} On 23 October 2026 we searched arXiv, bioRxiv and journal indexes for
action-conditioned or gene-conditioned world models and perturbation-conditioned video models of
live-cell imaging published in 2025--2026, and verified every cited work against its DOI or
arXiv record. We found no model that conditions the latent dynamics of live-cell video on the
perturbed gene, and none that evaluates zero-shot transfer of such a model to unseen genes.

\section{Data}\label{sec:data}

\paragraph{Movies.} The MitoCheck primary screen (screen A) imaged each siRNA-transfected well
every 30 minutes for two days (93 frames) \citep{neumann2010mitocheck}. We stream the movies at
the pyramid level of $512\times672$ pixels from the BioImage Archive's object store and never store
raw frames. Plates are split into \NPlatesTrain{} training, \NPlatesVal{} validation and
\NPlatesTest{} test plates; genes are split into \NGenesTrain{} training, \NGenesVal{} validation and
\NGenesTest{} test genes, with a second split that holds out whole protein complexes. A first
stage embedded \NMoviesStageOne{} movies on \NPlatesStageOne{} plates; a second stage, run on
training and validation plates only before any amended model was trained, selected
\NWellsStageTwoSelected{} further wells (all controls, every siRNA with a reproducible
phenotype, and random library wells up to a per-plate budget) and embedded
\NWellsStageTwoEmbedded{} of them on \NPlatesStageTwoEmbedded{} plates. The cache used for
training and validation holds \NCacheWells{} wells on \NCachePlates{} plates.

\paragraph{Missing frames.} \NWellsMissingFrames{} selected wells could not be embedded because
frames were missing from the archive's conversion, including \NWholePlatesMissing{} whole plates;
every missing frame was re-requested three times before the well was excluded. \NWellsMissingOnlyLate{}
of these wells miss only frames that the primary horizons do not use; the exclusion rule was
nevertheless fixed and applied identically to the test and screen-B data before they were read
(Section~\ref{sec:prereg}).

\paragraph{Independent reagents.} Screen B re-tested genes with independent siRNAs; it is
evaluated once, after the freeze, outside the gates.

\paragraph{Gene encodings.} Each gene is represented by a frozen vector from one of five sources
that use no MitoCheck data: genome-wide Perturb-seq pseudobulk responses in K562 cells
\citep{replogle2022perturbseq}, JUMP Cell Painting profiles of CRISPR knockouts and of ORF
overexpression \citep{chandrasekaran2023jump}, a spectral embedding of the STRING network
\citep{szklarczyk2023string} and mean-pooled ESM-2 protein embeddings \citep{lin2023esm2}.
Coverage of validation and test genes differs by source (Table~\ref{tab:coverage}).

\begin{table}[t]
\centering
\caption{Validation and test genes covered by each encoding source.}
\label{tab:coverage}
\begin{tabular}{lrr}
\toprule
source & validation genes & test genes \\
\midrule
Perturb-seq (K562) & \CovPerturbseqVal & \CovPerturbseqTest \\
JUMP CRISPR & \CovJumpCrisprVal & \CovJumpCrisprTest \\
JUMP ORF & \CovJumpOrfVal & \CovJumpOrfTest \\
STRING & \CovStringVal & \CovStringTest \\
ESM-2 & \CovEsmTwoVal & \CovEsmTwoTest \\
\bottomrule
\end{tabular}
\end{table}

\section{Model}\label{sec:model}

\paragraph{Latent states.} Each movie is cut into four $256\times256$ tiles and into 14 clips of
16 frames (8 hours); clips start every 6 frames (3 hours) for the first 13 clips. Each clip of
each tile is encoded by the frozen V-JEPA~2 ViT-L video encoder \citep{assran2025vjepa2}, its
tokens are averaged over time and pooled to $4\times4$ spatial tokens, and these are averaged to
one token per tile. A well's state at step $t$ is the set of its four tile tokens. DINOv2
\citep{oquab2023dinov2} and a video JEPA trained on live-cell microscopy serve as alternative
encoders.

\paragraph{Normalisation.} For each plate, clip and latent dimension, latents are centred and
scaled by the mean and standard deviation of that plate's negative-control wells
(\NNegPerPlateMin--\NNegPerPlateMax{} per plate), with the standard deviation floored at its 5th
percentile over plates. A global standardisation is kept as a sensitivity analysis.

\paragraph{Predictor.} A block-causal transformer takes the tile tokens of steps $\le t$, an
action token and a reagent (siRNA) token, and predicts the tile tokens at $t+1$ as a residual. The
action token is a learned per-gene embedding (seen genes), an MLP adapter of a frozen gene
encoding (zero-shot), or a constant (action-agnostic). Tile tokens carry a step embedding but no
position embedding, since fields of view are exchangeable. Two sizes are used: width
\WidthSmall{} and depth \DepthSmall{} (\ParamsSmall{}\,M core parameters) and width \WidthLarge{}
and depth \DepthLarge{} (\ParamsLarge{}\,M).

\paragraph{Training.} The loss is the L1 error of teacher-forced one-step predictions plus that
of a rollout of up to \RolloutSteps{} steps from a random context, as in V-JEPA~2-AC. We use
AdamW with learning rate \LearningRate{} and a one-cycle schedule over at most \MaxSteps{}
steps, batches of \BatchWells{} wells and siRNA-token dropout \SirnaDropout{}. The checkpoint with
the lowest validation primary error (evaluated every \EvalEvery{} steps) is kept, and training
stops \Patience{} steps after the last improvement.

\paragraph{Planner.} Given a start state and a target, the planner rolls the predictor out for
every candidate gene's encoding and ranks genes by the latent distance of the rolled-out state to
the target.

\section{Evaluation design}\label{sec:eval}

\paragraph{Primary horizons and wells.} The primary metric is the mean squared error of the
forecast latent at 24 and 36 hours after the first clip, in the normalised space. Most
knockdowns have no visible phenotype, so the primary wells for prediction are held-out-plate
wells whose siRNA carries a reproducible MitoCheck phenotype on training plates
(\NPrimaryWellsVal{} validation and \NPrimaryWellsTest{} test wells, frozen before training); the
primary genes for transfer are genes with a reproducible phenotype annotation
(\NPrimaryGenesVal{} validation and \NPrimaryGenesTest{} test genes). All wells and genes are
reported as secondary analyses.

\paragraph{Prediction (H1).} On primary wells of held-out plates, the gene-conditioned
predictor must beat the action-agnostic predictor and the replicate-drift baseline
(Section~\ref{sec:baselines}), each with a one-sided sign-flip test over plates at $p<0.05$, with
errors averaged over three seeds before testing.

\paragraph{Zero-shot transfer (H2).} On primary genes of the held-out gene split, at least one
encoding source must beat a model trained with the same encodings randomly permuted across genes
and the action-agnostic predictor, each with a 95\% bootstrap confidence interval over genes
below zero.

\paragraph{Planning (H3).} Phenotype planning ranks unseen genes for their ability to reach the
mean state of training wells with a MitoCheck phenotype (mitotic delay, polylobed nuclei,
binucleation); average precision must beat a label-permutation null (Benjamini--Hochberg
$q<0.05$) for at least two of the three phenotypes and a guilt-by-association baseline for at
least one. Unseen-gene truth sets hold \NTruthMitoticDelayVal, \NTruthPolylobedVal{} and
\NTruthBinuclearVal{} validation genes and \NTruthMitoticDelayTest, \NTruthPolylobedTest{} and
\NTruthBinuclearTest{} test genes. Trajectory planning ranks the gene that produced a held-out
trajectory; hit@100 must exceed random (binomial $p<0.05$) and endpoint retrieval with endpoints
imputed from the nearest training genes in the encoding space.

\section{Controls and baselines}\label{sec:baselines}

\paragraph{Baselines.} Copy-last; the time-slot mean; the plate's negative-control drift; the
action-agnostic predictor; \emph{replicate drift}, which adds to a well's own first state the
mean change of the same siRNA's wells on training plates (falling back to the gene's other
siRNAs, then to the plate's control drift); \emph{replicate endpoint}, the mean training-plate
latent of the same siRNA; a nearest-gene lookup in each encoding space; and, for planning,
guilt-by-association and endpoint retrieval. Replicate endpoints of unseen genes use those genes'
own observed wells, which the planner never sees, so they are reported as an oracle ceiling and
not as a gate comparator.

\paragraph{Controls.} Models trained on shuffled encodings (three seeds per source); an
evaluation-time shuffle; a static-clip control in which every clip repeats one frame; a context
swap in which the first state comes from another well with the same gene; a phenotype probe
applied to observed and forecast latents; probes that decode the plate from the latent within
replicate series; and a duplicate-movie check between splits by image fingerprints.

\section{Pre-registration, amendment and errata}\label{sec:prereg}

The design, hypotheses, splits and gates were committed before any test data were read. After a
first validation round on \NMoviesStageOne{} movies, the design was amended once, before any test
data were read, to add the per-plate normalisation, the primary horizons and wells, the
replicate baselines, three seeds and a fixed configuration grid (two predictor sizes, dropout
\Dropouts{}, weight decay \WeightDecays{}; \NGridRuns{} runs across \NFamilies{} model families),
a second stage of training data and a planning ground truth with more positives. The amendment
allows one round only: if the amended gates fail on validation, the models are frozen as they
are and the test is run once.

One erratum was recorded before any gate was computed: the amended trajectory-planning rule
compared the planner with replicate-endpoint retrieval, which for unseen genes uses the genes'
own observed wells and therefore cannot test zero-shot planning; the comparator was replaced by
nearest-gene-imputed endpoint retrieval, and the original baseline is reported as an oracle
ceiling. The rule for excluding wells with missing frames was likewise fixed for the test and
screen-B data before they were embedded. Configuration selection uses validation data only, the
validation gates were recorded before the freeze, and every number in Section~\ref{sec:results}
comes from models frozen before the test data were read. \TODO{author: review this section.}

\section{Results}\label{sec:results}

We report the pre-registered gates first, exactly as recorded before and after the freeze, and
then a set of robustness controls that we designed after seeing the test results. The controls
and their interpretation rules were written down and committed before any of them was run; every
number from them is post hoc and exploratory.

\subsection{Pre-registered gates}

\begin{table}[t]
\centering
\caption{Pre-registered decision gates on the validation and test splits. The test split was
evaluated once, after the freeze.}
\label{tab:gates}
\begin{tabular}{lll}
\toprule
Gate & Validation & Test \\
\midrule
Prediction (H1) & \SignalValPass & \SignalTestPass \\
Zero-shot transfer (H2) & \ZeroShotValPass & \ZeroShotTestPass \\
Planning (H3) & \PlanValPass & \PlanTestPass \\
\bottomrule
\end{tabular}
\end{table}

\paragraph{Prediction.} On the \HOneTestWells{} primary test wells on \SignalTestPlates{} plates,
the gene-conditioned predictor had \SignalTestAg\% lower error at 24 and 36 hours than the
action-agnostic predictor ($p=\SignalTestAgP$) and \SignalTestDrift\% lower error than replicate
drift ($p=\SignalTestDriftP$), so the prediction gate passes on test, as it did on validation
(\HOneValWells{} wells on \SignalValPlates{} plates; \SignalValAg\% and \SignalValDrift\% lower,
$p=\SignalValAgP$ and $p=\SignalValDriftP$; Table~\ref{tab:gates}). The comparison with the
action-agnostic predictor gave the same direction under the global preprocessing (validation
\HOneValGlobalAg\% \HOneValGlobalAgDir, $p=\HOneValGlobalAgP$; test \HOneTestGlobalAg\%
\HOneTestGlobalAgDir, $p=\HOneTestGlobalAgP$) and with DINOv2 latents (validation
\HOneValDinoAg\% \HOneValDinoAgDir, $p=\HOneValDinoAgP$; test \HOneTestDinoAg\%
\HOneTestDinoAgDir, $p=\HOneTestDinoAgP$). With TiJepa-1 latents the error was
\HOneValTijepaAg\% \HOneValTijepaAgDir{} on validation ($p=\HOneValTijepaAgP$) and
\HOneTestTijepaAg\% \HOneTestTijepaAgDir{} on test ($p=\HOneTestTijepaAgP$). After an
exploratory B-score correction of spatial plate effects, run after the test, the
gene-conditioned predictor still had \EOneValAg\% (validation, $p=\EOneValAgP$) and
\EOneTestAg\% (test, $p=\EOneTestAgP$) lower error than the action-agnostic predictor.

\paragraph{Screen B.} On screen B, a second MitoCheck screen with independent siRNAs whose
reagent token is unknown to the model, evaluated under the global preprocessing outside the
gates, the gene-conditioned predictor did not beat the action-agnostic predictor (difference
\ScreenBAg\% over \ScreenBPlates{} plates, $p=\ScreenBAgP$), although it beat replicate drift by
\ScreenBDrift\% ($p=\ScreenBDriftP$). The prediction rule therefore fails on screen B.

\paragraph{Zero-shot transfer.} The zero-shot gate passes the pre-registered rule on test for
one of five sources (string) with a CI bound at $\ZeroShotTestStringShufHi$; the rule has no
correction across sources; the same source fails on validation. We do not claim zero-shot
transfer. For that source, the trained model had a mean difference of
$\ZeroShotTestStringShuf$ (95\% CI $[\ZeroShotTestStringShufLo, \ZeroShotTestStringShufHi]$)
from the model trained on shuffled encodings and $\ZeroShotTestStringAg$
($[\ZeroShotTestStringAgLo, \ZeroShotTestStringAgHi]$) from the action-agnostic predictor, over
\ZeroShotTestStringGenes{} primary test genes.

\paragraph{Planning.} The planning gate fails on validation and test; planning is not
distinguishable from kNN-imputed retrieval. The difference in hit@100 between the planner and
nearest-gene-imputed endpoint retrieval was $\PlanValDiff$ on validation (95\% CI
$[\PlanValLo, \PlanValHi]$) and $\PlanTestDiff$ on test ($[\PlanTestLo, \PlanTestHi]$); the
oracle ceiling, which uses the unseen genes' own wells, reached hit@100 of \PlanValOracle{} and
\PlanTestOracle.

\subsection{Robustness controls (post hoc)}

All controls use the primary wells, the 24- and 36-hour metric, three seeds and the same
sign-flip test over plates as the prediction gate. Numbers are given for the test split, with
validation in parentheses.

\paragraph{Trained-shuffled control.} The gene-conditioned and action-agnostic predictors were
selected with different regularisation, so we trained the gene-conditioned model again at its
selected configuration with its action labels permuted: each training siRNA was mapped to
another, and every well received the permuted siRNA and that siRNA's gene, identically in
training and evaluation. This keeps the number of actions and the consistency between
replicates and removes the biological correspondence. The real model's error was
\RobOneRelTest\% \RobOneDirTest{} than the shuffled model's ($p=\RobOnePTest$; validation
\RobOneRelVal\% \RobOneDirVal, $p=\RobOnePVal$). The shuffled model had \RobShufAgTest\% lower
error than the action-agnostic predictor ($p=\RobShufAgPTest$; validation \RobShufAgVal\%,
$p=\RobShufAgPVal$), \RobShufFracTest\% (\RobShufFracVal\%) of the real model's gain. An
action-agnostic predictor trained at the gene-conditioned model's configuration differed from
the selected one by \RobNoneAgTest\% ($p=\RobNoneAgPTest$; validation \RobNoneAgVal\%,
$p=\RobNoneAgPVal$), so the regularisation difference does not explain the gain.

\paragraph{Additive memorisation.} A baseline that adds to the action-agnostic forecast a
fraction $\lambda$ of the average residual of the same siRNA's training-plate wells, with
$\lambda=\RobTwoLambda$ chosen on validation, had \RobTwoAgTest\% lower error than the
action-agnostic predictor ($p=\RobTwoAgPTest$; validation \RobTwoAgVal\%, $p=\RobTwoAgPVal$).
The gene-conditioned model's error was \RobTwoRelTest\% \RobTwoDirTest{} than this baseline's
(validation \RobTwoRelVal\% \RobTwoDirVal). The residuals are in-sample for the action-agnostic
model on training plates.

\paragraph{Reagent dependence.} With the siRNA token set to unknown at evaluation, the gain over
the action-agnostic predictor kept \RobThreeFracTest\% of its size ($p=\RobThreePTest$;
validation \RobThreeFracVal\%, $p=\RobThreePVal$), and \RobThreeGlobalFracTest\% under the
global preprocessing ($p=\RobThreeGlobalPTest$; validation \RobThreeGlobalFracVal\%,
$p=\RobThreeGlobalPVal$). Permuting gene and siRNA labels across wells within each plate made
the error \RobThreePermRelTest\% \RobThreePermDirTest{} than the action-agnostic predictor's
(validation \RobThreePermRelVal\% \RobThreePermDirVal).

\paragraph{Where the gain sits.} The gain over the action-agnostic predictor was absent at 3
hours (\RobFourThreehTest\%, $p=\RobFourThreehPTest$) and 12 hours (\RobFourTwelvehTest\%,
$p=\RobFourTwelvehPTest$) and appeared at 24 hours (\RobFourTwentyFourhTest\%,
$p=\RobFourTwentyFourhPTest$) and 36 hours (\RobFourThirtySixhTest\%,
$p=\RobFourThirtySixhPTest$). On non-primary library wells it was \RobFourNonPrimRelTest\%
($p=\RobFourNonPrimPTest$).

\paragraph{Interpretation rules.} The rules fixed before these controls were run require, for
any claim that the action input carries gene information, that the real model beat the
trained-shuffled model ($p<0.05$); this \RobIOneTest{} on test (\RobIOneVal{} on validation).
The rule for forecasts beyond additive memorisation \RobITwoTest, and the rule for a gene-level
rather than reagent-level gain \RobIThreeTest, because both require the first.

\section{Discussion}\label{sec:discussion}

The pre-registered prediction gate passed on validation and on test, and we report it as
recorded. A trained-shuffled control matches it: a model given a consistent but biologically
meaningless relabelling of the reagents forecasts as well as the gene-conditioned model. We
therefore make no gene-conditioning claim, and this paper reports a negative result. The gain
over the action-agnostic predictor shows that a consistent per-reagent label helps the forecast
at 24 and 36 hours; it does not show that the model has learned what the knocked-down gene does.
Screen B, where the reagents are new to the model and the gain over the action-agnostic
predictor disappears, is consistent with that reading.

The other two questions have negative answers under their own pre-registered rules. Zero-shot
transfer passes the pre-registered rule on test for one of five sources (string) with a CI bound
at $\ZeroShotTestStringShufHi$; the rule has no correction across sources; the same source fails
on validation. We do not claim zero-shot transfer. Planning fails on validation and test;
planning is not distinguishable from kNN-imputed retrieval.

Three lessons carry over to other perturbation world models. First, beating an action-agnostic
model is not evidence of action understanding when the two differ in what they can memorise; a
trained-shuffled control with the same label structure is the comparison that isolates
biological content. Second, replicate drift was a weak memorisation baseline here, and an
additive per-reagent correction of the action-agnostic forecast did at least as well as the
world model. Third, an independent reagent set is the most direct test of whether a learned
effect is gene-level. The negative result is specific to this model, these frozen video latents
and two days of low-resolution MitoCheck movies; richer encoders, longer horizons or multimodal
gene representations may change it, and the evaluation design can be reused to test them.

\section*{Licences and data availability}

MitoCheck movies and annotations (Image Data Resource idr0013; BioImage Archive S-BIAD865) are
CC0; the Perturb-seq pseudobulk data and STRING are CC~BY~4.0; JUMP Cell Painting profiles are
CC0. Model weights: V-JEPA~2 (MIT), DINOv2 (Apache-2.0), ESM-2 (MIT).
\TODO{author: code and data availability statement.}

\section*{Use of AI assistance}

Code, analyses and an initial manuscript draft were produced with the assistance of an AI coding
agent (Claude, Anthropic) under the author's direction; the author reviewed and is responsible
for all content. \TODO{author: review.}
